\section{Introduction}
\label{sec:intro}

Many of the decisions that software makes with a language model are small, typed and
repeated. Which team should receive a support ticket? Does a retrieved passage answer the
query? Does a cited source support, contradict or say nothing about a claim? Does a coding
agent's pending shell command go beyond what the user asked for? Does a web page contain
instructions addressed to the model rather than evidence? Which label of a caller-defined
taxonomy fits a document? In each case the application defines the admissible answers and
keeps the policy that acts on the answer (execute, hold, escalate) in its own code. What it
needs from the model is a probability distribution over those answers, returned quickly and
in a shape the code can check.

The common way to obtain such a decision from a large language model is to prompt it to
generate JSON. That route relies on the prompt to keep the output in shape, pays for
autoregressive decoding, and exposes probabilities only indirectly. \citet{kahneman2011}
describes fast, automatic judgments as ``System~1'' thinking; TypeSafe uses the name System
One for a hosted API built around this kind of judgment, served by its Jev model
\citep{typesafe_blog,typesafe_api}. Its public HTTP contract defines three typed primitives:
Choice (one of the caller's named candidates), Noul (a probability that a statement is true)
and Score (a position on an ordered rubric).

This paper describes Bobcat, a typed-decision model and serving contract built on the same
three primitives. Bobcat does not generate text for a decision. For each question it prefills a
compiled sequence (sharing the state prefix across questions where possible), reads the logits of the offered candidate identifiers at
the first answer position, and returns numbers; the host then builds the reply from the
caller's own candidate names and rubric, and an independent validator refuses any reply
outside a closed field set. The deployed model is Qwen3.8-27B \citep{qwen38_27b} with a
LoRA adapter \citep{hu2021lora} trained on labelled decisions, followed by temperature
calibration \citep{guo2017calibration}.

\paragraph{Relation to TypeSafe and Jev.}
We use TypeSafe's public HTTP documentation \citep{typesafe_api}, its MIT-licensed Python
SDK \citep{typesafe_sdk} and its public confidence adapter \citep{typesafe_adapter} as the
reference for the request and response shape and for the definition of the reported
\texttt{confidence} value. No output of Jev enters any training,
distillation, reward or calibration data in this work, and we never called Jev. The only Jev
output we use is the set of answers TypeSafe itself published for its workflow evaluations
\citep{typesafe_evals}, as a comparison row (Section~\ref{sec:jev}). Bobcat does not reproduce Jev's decisions, logits,
tokenizer, token counter or internal architecture; a third-party analysis of that
architecture \citep{archer_jev} informed some of our diagnostics but none of our claims.

\paragraph{What we measured.}
We built a six-task evaluation from native Korean sources with component-level splits,
counterfactual groups and a sealed final split (Section~\ref{sec:data}). We compared eight
pretrained backbones zero-shot through the same first-position readout
(Section~\ref{sec:leaderboard}), trained LoRA arms with cross-entropy, Brier and
proper-score REINFORCE objectives and a residual pointer head (Section~\ref{sec:arms}),
froze one candidate, \expone, and evaluated it once on the sealed final
(Section~\ref{sec:final}). We then stressed the served model for structural and numerical
fixedness and for semantic robustness (Section~\ref{sec:robust}), measured single-GPU
latency and throughput (Section~\ref{sec:serving}), ran public Korean and English benchmarks
(Section~\ref{sec:kbench}), and compared the served model with Jev on the workflow examples
TypeSafe published (Section~\ref{sec:jev}).

\paragraph{Contributions.}
\begin{itemize}
  \item A closed output contract for typed decisions in which the only model output is a
    vector of candidate logits, together with a compiler that keeps untrusted text out of
    the tokenizer's control vocabulary, assigns verified single-token identifiers, and
    encodes candidates piecewise so their positions are exact (Sections~\ref{sec:interface}
    and~\ref{sec:method}).
  \item A six-task evaluation of 6,400 decisions in 2,328 components, split by component,
    with 2,193 counterfactual groups and a final split that opens only for a frozen release
    manifest; and an account of an ordering bug in its first version
    (Section~\ref{sec:data}).
  \item Evidence that one epoch of LoRA supervised fine-tuning on a 27.8B hybrid backbone
    raises sealed-final task-macro accuracy by +10.2 points [+7.7, +12.3] over the same
    model zero-shot, including on a task that was held out of training, and two
    reinforcement-learning experiments showing that a proper-score REINFORCE objective,
    whose expected gradient equals the half Brier gradient (Proposition~1), matches but does
    not exceed direct Brier training, while correctness-only reward worsens calibration
    (Section~\ref{sec:experiments}).
  \item A question-level comparison with Jev on the 20 workflow examples TypeSafe published,
    scored with one code path for all models and checked for case-selection bias: parity in
    agreement with TypeSafe's reference, more probability on the reference answer, and a
    3 to 19 times slower deployment (Section~\ref{sec:jev}).
  \item Measurements of output fixedness (no structural violation in 12,846 requests,
    bit-identical repeats, batch and shared-prefix agreement, permutation and distractor
    stability) alongside a clear negative result: training made the model more willing to
    follow a wrong answer named inside the state (Section~\ref{sec:robust}).
  \item Single-GPU serving measurements for four workload profiles, a shared-prefix path that
    branches both the attention KV cache and the Gated DeltaNet convolution and recurrent
    state, and a vLLM engine study of the served artifact (Sections~\ref{sec:serving}
    and~\ref{sec:engine}).
\end{itemize}

\paragraph{Scope.}
The six-task evaluation is Korean only; English appears in a small author-written probe set
and in the 329 questions of the Jev comparison. None of the evaluation data has been reviewed by a person. The project's designated
quality reference, the direct-logit readout of GLM-5.3-Flash \citep{glm53flash}, has not
been run on the same evaluation, so the release gates defined against it are not judged
here. The model is larger than the 1 to 8B
parameter target of our own specification. Section~\ref{sec:limits} lists these and other
limitations together with the negative results.

\subsection{Related work}
\label{sec:related}

\paragraph{Reading decisions from logits.}
Scoring a fixed set of answer tokens at one position instead of sampling text is a
standard way to evaluate multiple-choice benchmarks such as MMLU \citep{hendrycks2021mmlu}
and HellaSwag \citep{zellers2019hellaswag}. Bobcat applies the same readout to caller-defined
candidates at serving time and treats the restriction to offered identifiers as a security
boundary rather than an evaluation convenience.

\paragraph{Calibration and proper scoring rules.}
The Brier score \citep{brier1950} and the log score are proper scoring rules. Post-hoc
temperature scaling \citep{guo2017calibration} fits one scalar on held-out data and leaves
the argmax unchanged. We use it as the only calibration step and compare a sampled
REINFORCE estimator \citep{williams1992} whose expected gradient equals that of the half
Brier score against direct Brier and cross-entropy training.

\paragraph{Efficient adaptation and serving.}
LoRA \citep{hu2021lora} trains low-rank updates of frozen weights. Gated DeltaNet
\citep{yang2024gdn} is a linear-attention layer with a gated delta-rule recurrence;
hybrids that interleave such layers with full attention carry a recurrent and a short
convolution state in addition to a KV cache, which matters when a shared prefix is reused
across branches. Prefix caching in serving engines such as vLLM \citep{kwon2023vllm}
reuses the KV blocks of a shared prompt; Section~\ref{sec:engine} measures such an engine on
our workloads.

\paragraph{Korean evaluation.}
Our product evaluation is built from KLUE machine reading comprehension and Wizard of Seoul
dialogues \citep{park2021klue}. Section~\ref{sec:kbench} reports public Korean
benchmarks: KoBEST \citep{kim2022kobest}, KMMLU \citep{son2024kmmlu}, HAE-RAE Bench
\citep{son2023haerae} and CLIcK \citep{kim2024click}.
